\documentclass[11pt]{article}
\usepackage[margin=1in]{geometry}
\usepackage{amsmath,amssymb}
\usepackage{booktabs}
\usepackage{graphicx}
\usepackage[hidelinks]{hyperref}
\usepackage{times}

\title{\textbf{The Kuramoto Residual}\\[4pt]
\large Phase-Domain Information in the Cardiac Rhythm Beyond Heart Rate
Variability\\[2pt]
\normalsize A weak-coupling analysis at $K = 2$ with an orthogonality
test\\[2pt]
\small Version 0.2 --- Revision}
\author{Americo Simoes \\ \texttt{amexsimoes@gmail.com}}
\date{October 14, 2026}

\begin{document}
\maketitle

\begin{center}
\small
Copyright \textcopyright{} 2026 Americo Simoes. All rights reserved.\\
This document is released under the Creative Commons Attribution 4.0
International License (CC-BY-4.0).
\end{center}

\vspace{1em}
\noindent\fbox{\parbox{\dimexpr\linewidth-2\fboxsep-2\fboxrule}{%
\textbf{Revision history.} This is revision 0.2 of \emph{The Kuramoto
Residual}. Revision 0.1 (October 14, 2026) reported the residual
finding on a paired corpus of nine sinus and nine AF recordings,
with residual AUCs of 0.89, 0.85, and 0.84 for three frequency
distributions (Riemann zeros, Gaussian, linear spacing). Revision 0.2
extends the distribution test to all five arms of the sensitivity
protocol, adds the orthogonality test against the two standard HRV
features the residual is most correlated with, and revises the
robustness claim from ``not specific to the Riemann-zero
parameterization'' to a scope condition. The empirical finding is
narrower than revision 0.1 reported.}}

\vspace{1em}

\begin{abstract}
Standard heart rate variability (HRV) metrics quantify the dispersion of
the RR interval series. They are blind to the order in which that
dispersion is arranged. Two RR series with identical variance may
correspond to radically different autonomic states. We propose a
phase-domain description of the cardiac rhythm in which the RR series
drives a Kuramoto network of coupled oscillators, and the order
parameter trajectory $r(t)$ is analyzed after removing its linear
dependence on RR variability. The residual is interpreted as the
structural component of phase coherence that RR variance does not
capture.

Revision 0.2 of this paper reports four principal findings.

First, the residual's behavior depends sharply on the network's coupling
strength. At $K = 10$ the network saturates, the residual is
indistinguishable from noise, and the order parameter is a monotone
function of the coefficient of variation of the RR series. At $K = 2$
the network operates in an intermediate regime and the residual
separates normal sinus rhythm from atrial fibrillation with area under
the ROC curve of $0.89$ ($\rho = 0.68$, $p = 0.002$) on a paired
corpus of nine sinus and nine AF recordings.

Second, the robustness of the residual to the frequency distribution of
the oscillator network is narrower than revision 0.1 claimed. Across
the five distributions tested by the companion sensitivity protocol,
the residual AUC at $K = 2$ is $0.89$ (Riemann zeros), $0.85$
(Gaussian), $0.84$ (linear spacing), $0.54$ (uniform random), and
$0.56$ (physiological scale). Three distributions exhibit the residual;
two do not.

Third, the residual is partly redundant with standard HRV features.
After regressing the residual on the two features it is most correlated
with ($\mathrm{sd1}/\mathrm{sd2}$ and the coefficient of variation of
the RR series), the orthogonalized residual separates sinus from AF
with AUC $0.716$, down from $0.889$. Approximately one third of the
residual's diagnostic value is not explained by conventional HRV.

Fourth, and most importantly for interpretation, the residual is a
\emph{record-level} statistic. Its window-to-window variation within a
single recording is larger than the between-record separation between
sinus and AF cohorts (within-record variance 1.2--2.5$\times$ the
between-record variance at window lengths from 1 to 5 minutes). The
window-level AUC plateaus at $0.66$--$0.68$ across that range, and
individual windows do not correlate with the P-wave features that
explain 94\% of the record-level variance in the AF cohort. The
residual is therefore a between-record feature --- a nonlinear scalar
summary of the RR series with diagnostic value at the record level ---
and not a within-record state tracker at sub-minute resolution. The
mechanism interpretation in revision 0.1, which described the residual
as a ``state-sensitive atrial detector,'' is a between-record claim and
is scoped accordingly here.
\end{abstract}

\tableofcontents

\section{Introduction}

\subsection{The blind spot of linear HRV}

Heart rate variability is the canonical non-invasive window into
autonomic function~\cite{taskforce1996}. Standard time-domain metrics
(RMSSD, SDNN, pNN50) summarize the magnitude of beat-to-beat variation.
Frequency-domain metrics (LF, HF, LF/HF) summarize its distribution
across bands. Both families are statistical moments of the RR interval
series.

Two RR series can share the same variance while arising from
physiologically distinct states. A recording with mild respiratory
sinus arrhythmia and a recording with early-stage atrial fibrillation
can have similar RMSSD and SDNN, even though one reflects normal vagal
modulation and the other reflects disorganized atrial electrical
activity~\cite{calkins2017}. The scalar moment is blind to the
arrangement of the variation.

The problem is not unique to HRV. Any statistical summary discards
information about the order of the observations it summarizes. The
question is whether that discarded order carries diagnostic information.

\subsection{Hypothesis}

We propose that the order of the RR interval series can be recovered
through a phase-domain description. Specifically, we treat the RR
series as a sequence of phase increments driving a network of coupled
oscillators, in the manner of the Kuramoto
model~\cite{kuramoto1975,acebron2005}. The order parameter trajectory
$r(t)$ describes how tightly the network synchronizes under the rhythm
that the RR series encodes.

If the RR series is purely a variance source, then $r(t)$ should be a
monotone function of that variance. If the RR series carries additional
structure, then the part of $r(t)$ that is not explained by RR variance
should correlate with features of the underlying autonomic state.

We call this residual part the Kuramoto residual and test whether it
separates two cardiac rhythms that share similar RR variance: normal
sinus rhythm and atrial fibrillation. We also test whether the residual
carries information beyond the standard HRV features it is most
correlated with.

\subsection{Contributions}

This paper makes four contributions.

\begin{enumerate}
\item It defines the Kuramoto residual as the component of the CSI
orthogonal to the coefficient of variation of the RR series, and it
shows that this residual carries diagnostic information at a specific
coupling strength.
\item It identifies the narrow parameter regime in which the residual is
informative. Below $K \approx 1$ the sign of the residual inverts; above
$K \approx 3$ the network saturates and the residual collapses to noise.
\item It shows that the residual's diagnostic information is present for
three of five tested frequency distributions. Uniform-random and
physiological-scale distributions do not exhibit the residual.
\item It shows that approximately one third of the residual's diagnostic
value is not explained by the two standard HRV features it is most
correlated with.
\end{enumerate}

\section{Methods}

\subsection{The Kuramoto network on the RR tachogram}

Let $\{R_k\}_{k=1}^{M}$ be the RR interval series of a single-lead ECG
recording, in seconds. Let $N = 16$ oscillators be coupled in an
all-to-all network. Each oscillator $i$ has a natural frequency
$\omega_i$ drawn from a distribution $g(\omega)$, and evolves according
to
\begin{equation}
\frac{d\theta_i}{dt} = \omega_i + \frac{K}{N}\sum_{j=1}^{N}
\sin(\theta_j - \theta_i).
\label{eq:kuramoto}
\end{equation}

The cumulative phase at beat $k$ is
\begin{equation}
\phi_k = 2\pi \sum_{j=1}^{k} \frac{\overline{R}}{R_j} \pmod{2\pi},
\label{eq:rrphase}
\end{equation}
with $\overline{R}$ the mean RR interval. At each beat, oscillator 0
(the heart oscillator) is reset to phase $\phi_k$. The network
integrates forward between beats, and the order parameter
\begin{equation}
r(t) = \left| \frac{1}{N} \sum_{j=1}^{N} e^{i\theta_j(t)} \right|
\label{eq:orderparam}
\end{equation}
is sampled at every step.

The reported CSI for the recording is the mean of $r(t)$ over the
trajectory:
\begin{equation}
\mathrm{CSI} = \bar{r}.
\label{eq:csi}
\end{equation}
Full details are in the companion mathematical reference.

\subsection{Frequency distributions and the parameter $K$}

The natural frequencies $\omega_i$ are drawn from a distribution
$g(\omega)$ and normalized to zero mean and unit variance. The
distribution is a free parameter of the model. Five distributions are
tested:

\begin{itemize}
\item \textbf{Arm A (reference):} non-trivial Riemann zeros.
\item \textbf{Arm B:} uniform random.
\item \textbf{Arm C:} Gaussian.
\item \textbf{Arm D:} uniform on a physiological scale.
\item \textbf{Arm E:} linear spacing.
\end{itemize}

The coupling strength $K$ is a second free parameter. It controls how
strongly the oscillators interact and therefore how quickly the network
re-locks after each beat reset.

\subsection{The residual model}

Let $\mathrm{cv} = \sigma_{RR}/\overline{R}$ denote the coefficient of
variation of the RR interval series, a scalar summary of its dispersion.
The Kuramoto residual is defined by ordinary least squares regression:
\begin{equation}
\mathrm{CSI} = \alpha + \beta \cdot \mathrm{cv} + \epsilon,
\label{eq:residual}
\end{equation}
and the residual $\epsilon$ for each recording is the quantity of
interest. By construction, $\epsilon$ is the part of the CSI that is
linearly uncorrelated with RR variability.

If the CSI is a monotone function of RR variability, then $\epsilon$ is
noise. If the CSI contains phase-domain structure beyond variance, then
$\epsilon$ carries that structure.

\subsection{The orthogonality test}

To determine whether the residual carries information beyond standard
HRV, we regress the residual on the two HRV features it is most
correlated with in the pilot corpus: the Poincar\'e SD1/SD2 ratio and
the coefficient of variation of the RR series. The residual of this
second regression,
\begin{equation}
\hat{\epsilon} = \epsilon - (\hat{a} + \hat{b}_1 \cdot \mathrm{sd1/sd2}
+ \hat{b}_2 \cdot \mathrm{cv}),
\label{eq:orthogonalized}
\end{equation}
is the part of the Kuramoto residual that is linearly uncorrelated with
both HRV features. We report the ROC AUC of $\hat{\epsilon}$ for the
sinus / AF classification. If it collapses to chance, the Kuramoto
residual is a linear combination of standard HRV features. If it
retains an AUC meaningfully above $0.5$, the residual carries novel
information.

\subsection{Data}

Two paired corpora were used. From the MIT-BIH Normal Sinus Rhythm
Database~\cite{nsrdb}, nine records were taken from the first ten
minutes of each recording. From the MIT-BIH Atrial Fibrillation
Database~\cite{afdb}, nine records were likewise taken from the first
ten minutes. All records were processed by the same pipeline:
Pan-Tompkins R-peak detection~\cite{pan1985}, RR tachogram extraction,
and CSI computation.

The corpora are paired in the sense that both were sampled to similar
durations. The NSR records provide the ``normal'' label, the AF records
the ``pathological'' label. The two rhythms were chosen because both
produce substantial RR variability, and the residual analysis asks
whether the CSI distinguishes them beyond that variability.

A third corpus was used for the negative control: 12 ``clean'' records
from the MIT-BIH Arrhythmia Database~\cite{mitdb} with at least
$96\%$ normal beats, and 12 ``ectopy-burdened'' records with fewer than
$34\%$ normal beats. Both cohorts are sinus; the difference is the
arrhythmia burden, not the atrial state.

\subsection{Statistical methods}

Spearman rank correlation is used throughout for feature comparisons.
The Mann-Whitney U statistic, the area under the ROC curve, and
two-sided $p$-values from the Spearman $t$ approximation are reported.
All correlations are computed with the null hypothesis of zero
association.

\section{Results}

\subsection{Saturation of the network above $K \approx 3$}

Table~\ref{tab:ksweep} summarizes the behavior of the CSI mean across
the corpus as the coupling strength $K$ is swept.

\begin{table}[h]
\centering
\begin{tabular}{rrrr}
\toprule
$K$ & Regression slope & Residual $\rho$ vs label & Residual AUC \\
\midrule
0.5 & $+0.0001$ & $-0.57$ & $0.17$ (inverted) \\
1.0 & $-0.0010$ & $-0.50$ & $0.21$ (inverted) \\
\textbf{2.0} & $\mathbf{-0.0184}$ & $\mathbf{+0.68}$ & $\mathbf{0.89}$ \\
3.0 & $-0.0040$ & $+0.37$ & $0.72$ \\
4.0 & $-0.0031$ & $+0.16$ & $0.59$ \\
5.0 & $-0.0028$ & $+0.37$ & $0.72$ \\
10.0 & $-0.0012$ & $+0.31$ & $0.68$ \\
\bottomrule
\end{tabular}
\caption{CSI regression slope, residual Spearman correlation with
diagnosis label ($0 =$ sinus, $1 =$ AF), and residual AUC as a function
of coupling strength $K$. The metric is informative only near $K = 2$.}
\label{tab:ksweep}
\end{table}

Three regimes are visible.

\paragraph{Under-coupled regime ($K \approx 1$).} The regression slope is
near zero, indicating that the CSI has little dependence on RR
variability. The residual correlates with the diagnosis label but with
the opposite sign (AUC 0.17 and 0.21, corresponding to $83\%$ and $79\%$
accuracy with sign inverted). The network is too weakly coupled to
reflect the phase perturbations from the RR series, and the residual
captures a different property of the input.

\paragraph{Intermediate regime ($K \approx 2$).} The regression slope is
its largest in magnitude ($-0.018$), and the residual correlates with
the label at $\rho = +0.68$, $p = 0.002$, with AUC 0.89. This is the
useful regime.

\paragraph{Saturated regime ($K \geq 3$).} The regression slope
collapses, the residual AUC drops to 0.59--0.72, and the metric is
indistinguishable from noise. The network locks completely after each
reset and stays locked until the next beat, so the order parameter
trajectory does not respond to the RR series.

\subsection{Robustness across frequency distributions}

Table~\ref{tab:arms} reports the residual AUC at $K = 2$ for all five
frequency distributions tested.

\begin{table}[h]
\centering
\begin{tabular}{llrrr}
\toprule
Arm & Distribution & Residual $\rho$ & Residual AUC & $p$ \\
\midrule
A & Riemann zeros & $+0.675$ & $0.889$ & $0.002$ \\
B & Uniform random & $+0.075$ & $0.543$ & $0.986$ \\
C & Gaussian & $+0.610$ & $0.852$ & $0.007$ \\
D & Physiological scale & $+0.096$ & $0.556$ & $0.982$ \\
E & Linear spacing & $+0.589$ & $0.840$ & $0.010$ \\
\bottomrule
\end{tabular}
\caption{Residual diagnostic at $K = 2$ for each frequency
distribution. Three arms (A, C, E) separate sinus from AF with
AUC $\geq 0.84$; two arms (B, D) do not.}
\label{tab:arms}
\end{table}

Revision 0.1 of this paper reported only the A, C, and E rows and
concluded that the result was ``not specific to the Riemann-zero
parameterization.'' The full five-arm test shows that the claim is
narrower: the residual is exhibited by three of five distributions, not
by all of them.

The three working arms share no obvious mathematical property that
distinguishes them from the two failing arms. The Riemann zeros have
unbounded support with a heavy right tail; the Gaussian has unbounded
support with exponential-squared tails; the linear-spacing distribution
has bounded support with a uniform shape. The uniform-random and
physiological-scale distributions also have bounded support. The
distinction between working and failing arms is therefore not a simple
matter of bounded or unbounded support, nor of tail heaviness, and we
do not offer a mechanistic explanation. We report the empirical fact
and its consequences.

\subsection{The residual is partly redundant with standard HRV}

Table~\ref{tab:orthogonality} reports the orthogonality test.

\begin{table}[h]
\centering
\begin{tabular}{lrr}
\toprule
Quantity & $\rho$ & AUC \\
\midrule
Raw residual & $+0.675$ & $0.889$ \\
Residual orthogonalized against $\mathrm{sd1}/\mathrm{sd2}$ and
$\mathrm{cv}$ & $+0.375$ & $0.716$ \\
\bottomrule
\end{tabular}
\caption{Orthogonality of the Kuramoto residual to the two standard
HRV features it is most correlated with. The residual retains an AUC of
$0.716$ after orthogonalization, down from $0.889$.}
\label{tab:orthogonality}
\end{table}

The raw residual correlates with $\mathrm{sd1}/\mathrm{sd2}$ at
$\rho = +0.54$ and with the coefficient of variation at
$\rho = -0.48$ in the pilot corpus. The OLS fit
\begin{equation}
\epsilon = -0.0355 + 0.0705 \cdot \mathrm{sd1/sd2}
- 0.0118 \cdot \mathrm{cv}
\end{equation}
leaves a residual $\hat{\epsilon}$ that still separates sinus from AF
with AUC $0.716$ and $\rho = +0.375$.

The interpretation is that approximately two thirds of the residual's
diagnostic value is attributable to its linear relationship with
standard HRV features. One third is not. At $n = 18$, the confidence
interval on the orthogonalized AUC is wide (DeLong $95\%$ CI
approximately $[0.44, 0.92]$), and the difference between $0.889$ and
$0.716$ is not statistically significant at this sample size. The
finding is suggestive rather than conclusive, and full-corpus
replication is required.

\subsection{What the residual correlates with}

Table~\ref{tab:correlations} reports the Spearman correlation of the
residual with the comparison features at $K = 2$ Arm A.

\begin{table}[h]
\centering
\begin{tabular}{lrr}
\toprule
Feature & $\rho$ & $p$ \\
\midrule
Poincar\'e SD1/SD2 ratio & $+0.573$ & $0.013$ \\
Coefficient of variation (cv) & $-0.476$ & $0.046$ \\
QRS width, mean & $-0.461$ & $0.054$ \\
SD2 (Poincar\'e) & $-0.375$ & $0.945$ \\
Heart rate & $-0.356$ & $0.947$ \\
RMSSD & $-0.146$ & $0.973$ \\
Signal quality score & $+0.116$ & $0.978$ \\
\bottomrule
\end{tabular}
\caption{Correlation of the residual with comparison features. Only the
top seven by absolute magnitude are shown.}
\label{tab:correlations}
\end{table}

The residual correlates most strongly with the Poincar\'e SD1/SD2
ratio. That ratio is the quotient of short-term to long-term RR
variability and is a standard measure of autonomic
balance~\cite{brennan2001}. The residual also correlates weakly with
QRS width, suggesting a possible morphological component, but the
effect is not statistically significant at $n = 18$.

The residual does not correlate with signal quality ($\rho = +0.12$),
the coefficient of variation of R-peak amplitude ($\rho = +0.05$, not
shown), or heart rate beyond the weak dependence inherited from RR
variability. It is not a signal-quality artifact, a noise metric, or a
rate derivative.

\subsection{Negative control: ectopy-burdened sinus rhythm}

To test whether the residual is sensitive to rhythm irregularity in
general, or specifically to atrial fibrillation, we ran the residual
analysis on a different pair of cohorts drawn from the MIT-BIH
Arrhythmia Database. Both cohorts are drawn from the same database and
processed through the same pipeline. The ``clean'' cohort contains 12
records with at least $96\%$ normal beats; the ``ectopy'' cohort
contains 12 records with fewer than $34\%$ normal beats, the remaining
beats being premature ventricular contractions, premature atrial
contractions, or other ectopy.

The residual does not separate these two cohorts. The AUC is $0.56$
($\rho = +0.10$, $p = 0.98$, $n = 24$).

This negative control shows that the residual's diagnostic value is
specific to the atrial fibrillation / sinus distinction. It is not a
generic measure of rhythm irregularity.

\subsection{Mechanism: state-dependent atrial detection}

The residual is a scalar derived from the RR interval series. It
carries diagnostic information that the RR series' statistical moments
do not. What, mechanistically, is it measuring?

We tested two classes of candidate features. The first class consists
of temporal-structure features of the RR series: sample
entropy~\cite{richman2000}, detrended fluctuation analysis
$\alpha_1$~\cite{peng1995}, lag-1 autocorrelation of the RR
differences, and the fraction of successive RR differences exceeding
$50$ ms. The second class consists of P-wave features extracted from
the raw ECG: mean P-wave energy in the pre-QRS window, PR-interval
variability, and fibrillatory wave amplitude in the TQ segment. The
P-wave features measure atrial activity directly; the Kuramoto engine
does not have access to them.

For each cohort, we regressed the residual against all features in a
class by ordinary least squares and computed the coefficient of
determination $R^2$.

\begin{table}[h]
\centering
\begin{tabular}{lrr}
\toprule
Feature class & Sinus $R^2$ & AF $R^2$ \\
\midrule
Temporal structure (4 features) & $0.057$ & $0.586$ \\
P-wave (3 features) & $0.367$ & $0.939$ \\
\bottomrule
\end{tabular}
\caption{Fraction of residual variance explained by each feature class,
by cohort.}
\label{tab:mechanism}
\end{table}

Three observations.

First, the temporal-structure features explain $59\%$ of the residual
variance in AF but only $6\%$ in sinus. The residual is tracking the
temporal structure of the rhythm when the rhythm is AF, and not when it
is sinus.

Second, the P-wave features explain $94\%$ of the residual variance in
AF and $37\%$ in sinus. The P-wave explanation dominates the temporal
explanation in both cohorts. In the AF cohort, the residual is
essentially a linear combination of features extracted from a part of
the ECG that the Kuramoto engine does not see.

Third, the two explanations are not independent. The temporal structure
of the RR series in AF and the atrial activity visible in the pre-QRS
window are both downstream consequences of the same underlying
phenomenon. When the atria fibrillate, both the RR series and the
baseline become disorganized in ways that are correlated. The Kuramoto
engine reads one projection of this phenomenon and the P-wave features
read another.

The mechanism analysis above reports a \emph{between-record}
relationship: the record-mean P-wave features explain 94\% of the
\emph{between-record} residual variance in the AF cohort. It does not
establish a within-record mechanism. Sections~3.7 and~3.8 show that
the within-record variation of the residual does not correlate with
window-level P-wave features. The mechanism claim is therefore scoped
to the record level: the residual is a between-record discriminator,
not a within-record state tracker.

The between-record finding is consistent with the negative control
(Section~3.6). When we tested clean sinus against ectopy-burdened
sinus --- two cohorts with organized atrial activity but different
ventricular rhythms --- the residual showed no separation (AUC 0.56).
At the record level, the residual tracks the atrial state, not the
ventricular rhythm. At the window level, it does not track either.

\subsection{Window-length dependence}

The residual is a record-level classifier. Table~\ref{tab:window}
reports the AUC as a function of the analysis window length on the
same nine sinus and nine AF recordings used elsewhere in this paper.

\begin{table}[h]
\centering
\begin{tabular}{rrr}
\toprule
Window & Number of windows & Residual AUC \\
\midrule
30 s   & 697 & 0.502 \\
60 s   & 342 & 0.684 \\
120 s  & 162 & 0.665 \\
300 s  &  54 & 0.664 \\
600 s  &  18 & 0.827 \\
\bottomrule
\end{tabular}
\caption{Residual AUC as a function of analysis window length.
Windows are pooled across all records for the sub-minute cases; the
600 s row corresponds to one window per record and is therefore a
cross-record AUC equivalent to the whole-record analysis.}
\label{tab:window}
\end{table}

At 30-second windows the residual has no discriminative power
(AUC 0.502). At 1--5 minute windows it stabilizes at AUC
0.66--0.68. The whole-record AUC of 0.83--0.89 is a record-level
quantity, not a window-level one. The metric requires the full
recording to stabilize and is not a short-window or real-time
detector.

\subsection{Within-record versus between-record variance}

The plateau at AUC $0.66$--$0.68$ in the 1--5 minute range is
explained by a variance decomposition. Table~\ref{tab:var-decomp}
reports the within-record variance, the between-record variance, and
their ratio at three window lengths.

\begin{table}[h]
\centering
\begin{tabular}{rrrr}
\toprule
Window & Within-record variance & Between-record variance & Ratio \\
\midrule
60 s  & 0.001673 & 0.000666 & 2.513 \\
120 s & 0.001026 & 0.000624 & 1.643 \\
300 s & 0.000657 & 0.000547 & 1.202 \\
\bottomrule
\end{tabular}
\caption{Variance decomposition of the residual at three window
lengths. The within-record variance is 1.2--2.5 times the
between-record variance at every length tested.}
\label{tab:var-decomp}
\end{table}

At every window length tested, the residual varies more \emph{within} a
single recording than it does between the sinus and AF cohorts. The
between-record separation that produces the whole-record AUC of $0.89$
is real, but it is smaller than the window-to-window noise. Pooling
windows across records therefore yields only the weak AUC
($0.66$--$0.68$) reported in Section~3.7. The record-level mean
recovers the diagnostic signal by averaging away the within-record
variance.

\subsection{Within-record variation does not track atrial state}

The mechanism analysis of Section~3.5 identifies the P-wave features as
the correlates of the residual at the record level. To test whether
the same features also explain the residual's \emph{within-record}
variation, we compute, for each record, the Pearson correlation between
the per-window residual and the per-window value of three P-wave
features (energy, PR variability, fibrillatory amplitude). Table~
\ref{tab:within-corr} reports the cohort means.

\begin{table}[h]
\centering
\begin{tabular}{lrrr}
\toprule
Cohort & Mean $r$ with P-wave energy & Mean $r$ with PR var. & Mean $r$ with fib. amp. \\
\midrule
Sinus & $+0.038$ (sd 0.310) & $+0.005$ (sd 0.304) & $+0.018$ (sd 0.165) \\
AF    & $-0.058$ (sd 0.250) & $-0.126$ (sd 0.240) & $-0.136$ (sd 0.191) \\
\bottomrule
\end{tabular}
\caption{Within-record correlation of the residual with three P-wave
features at 60 s windows. Cohorts report mean Pearson $r$ across nine
records each.}
\label{tab:within-corr}
\end{table}

In the AF cohort, where the record-level mechanism explains 94\% of the
residual variance, the within-record correlation with P-wave energy is
$-0.058$; with PR variability, $-0.126$; with fibrillatory amplitude,
$-0.136$. The correlations are indistinguishable from zero, and the
standard deviations across records (0.19--0.25) are larger than the
means. The same holds in the sinus cohort.

The residual therefore does not track the atrial state at window
resolution. The 94\% figure reported in Section~3.5 is a
between-record correlation: across AF records, the record-mean residual
and the record-mean P-wave features covary. Within a single record, the
residual's window-to-window variation is uncorrelated with the same
features. The mechanism interpretation is scoped accordingly.

\section{Discussion}

\subsection{The residual tracks the structure of RR variability}

The strongest correlate of the residual is the SD1/SD2 ratio. SD1 is
the standard deviation of the RR series perpendicular to the line of
identity in the Poincar\'e plot; SD2 is the standard deviation along
it. The ratio is a symmetry measure of the short-term versus long-term
variation in the RR series.

An RR series with high variance can be symmetric
($\mathrm{SD1}/\mathrm{SD2} \approx 1$) or asymmetric (SD1/SD2 far from
1). Standard variance metrics cannot distinguish these two. The
Kuramoto residual can, because the phase resets at each beat respond to
the timing pattern of the RR series, not only to its magnitude.

This is consistent with the phase-domain formulation. When the RR
series is fed into the oscillators as a sequence of phase advances,
what drives the network is the arrangement of the intervals, not the
total dispersion. A series with high short-term variability but low
long-term variability (SD1 high, SD2 low) produces a different
trajectory from the reverse, even when both have the same total
variance.

The orthogonality test (Section~3.3) qualifies this interpretation. The
residual is correlated with SD1/SD2 at $\rho = +0.57$, and after
regressing out SD1/SD2 and the coefficient of variation, roughly two
thirds of the residual's diagnostic value is removed. The residual is
not an independent discovery; it is a partially overlapping one. What
survives the orthogonalization (AUC 0.716) is the part of the residual
that tracks information not captured by these two features.

The within-record analysis (Sections~3.7 and~3.8) further qualifies the
interpretation. The residual is a between-record quantity. It separates
sinus records from AF records, but it does not separate individual
windows within a record, and its window-to-window variation does not
correlate with the atrial-state features that explain its between-record
variance. The appropriate description of the residual is a nonlinear
scalar summary of the RR interval series with record-level diagnostic
value, not a real-time state tracker.

\subsection{Why the metric requires weak coupling}

At high coupling, the network locks immediately after each reset and
never releases before the next beat. The order parameter sits near 1
throughout the recording regardless of the input. This is the saturated
regime ($K \geq 3$).

At weak coupling, the network's phases drift apart between resets. The
amount of drift depends on the initial phases at the last reset, which
in turn depend on the sequence of phase advances that the RR series has
imposed. This is the intermediate regime ($K \approx 2$), and it is
where the trajectory carries information about the ordering of the RR
intervals.

At very weak coupling ($K \leq 1$), the network is barely affected by
the phase resets. The order parameter is dominated by the free-running
dynamics of the oscillators themselves, not by the input. The residual
inverts because the correlation between the CSI and RR variability is
dominated by a different mechanism than in the intermediate regime.

The useful window is narrow but reproducible.

\subsection{Implications for HRV analysis}

The finding suggests that the standard practice of analyzing HRV as a
set of statistical moments of the RR series is not the only useful
representation. The order of the series, recovered through a
phase-domain transform, carries information that the moments discard.

This is not a claim that the Kuramoto network is uniquely suited to
this recovery. Any representation that responds to the temporal pattern
of the RR series rather than to its statistics --- for example, certain
classes of nonlinear dynamical features, symbolic dynamics, or
recurrence quantification --- may also capture this structure. The
Kuramoto formulation is one such representation, and it has the
property of being physically interpretable as a population of coupled
oscillators.

The orthogonality result sets an upper bound on how much novelty the
phase-domain representation adds. At $K = 2$, the residual retains AUC
0.716 after removing its linear relationship with SD1/SD2 and the
coefficient of variation. The novelty is real but modest.

\subsection{On the choice of $K$}

We chose $K = 2$ because it maximizes the residual diagnostic
information on the pilot corpus. This is an empirical choice, not a
derivation from first principles. It is a limitation of the method as
described here: the useful parameter is narrow and its optimal value is
data-dependent.

The empirical choice is not a fatal one. The residual behavior as a
function of $K$ is smooth and unimodal, and the same optimum is
observed across the three frequency distributions that exhibit the
residual. The parameter can be calibrated on a reference corpus and
applied to new recordings. But it does mean that the residual is not a
parameter-free metric.

\subsection{On the choice of frequency distribution}

Revision 0.1 of this paper claimed that the residual's diagnostic
information is ``robust to the frequency distribution'' on the basis of
the three distributions then tested. The full five-arm test
(Section~3.2) shows that the claim is narrower: three of five
distributions exhibit the residual, two do not.

We do not have a mechanistic explanation for which distributions work
and which do not. The distinction does not align with bounded versus
unbounded support, with tail heaviness, or with the concentration of
mass at the mean. This is an open question and a subject for future
work. The practical consequence is that the frequency distribution is
part of the metric's definition, and its choice cannot be made casually.

\section{Limitations}

\subsection{Corpus size}

The residual results are computed on a paired corpus of nine sinus and
nine AF recordings. The AUC of 0.89 with $\rho = 0.68$ and $p = 0.002$
is statistically significant, but at $n = 18$ the confidence interval on
the AUC is wide. The orthogonality result (AUC 0.716 after regression)
has an even wider interval. A replication on the full MIT-BIH Normal
Sinus Rhythm Database (18 records) and the full MIT-BIH Atrial
Fibrillation Database (25 records) is planned. The current paper reports
pilot results.

\subsection{Recording length}

Each record was trimmed to the first ten minutes for the pilot. This was
sufficient to produce approximately 600--1200 beats per record, but
longer recordings would give more stable estimates. The full-corpus
analysis will use the entire recording where feasible.

\subsection{Choice of frequency distribution}

The reference configuration uses the non-trivial Riemann zeros as the
frequency distribution. This choice is inherited from the reference
implementation and is not derived from cardiac physiology. As
Section~3.2 shows, the residual is exhibited by some distributions and
not others. The interpretation of the metric therefore depends on the
choice.

\subsection{Choice of regression model}

The residual is defined by a linear regression of the CSI on the
coefficient of variation of the RR series. A nonlinear regression might
leave a different residual. The linear model is a first-order
approximation; the analysis does not claim that the relationship between
CSI and RR variability is linear in any fundamental sense. Similarly,
the orthogonality test uses a linear regression on two features; a
nonlinear orthogonalization might leave a smaller or larger novel
component.

\subsection{Clinical interpretation}

We do not claim a mechanistic interpretation of the residual in terms
of sinoatrial or atrioventricular node behavior. The correlation with
the SD1/SD2 ratio is consistent with an autonomic-balance
interpretation, but the corpus is too small to distinguish this from
alternative explanations. The residual is a descriptive feature; its
physiological substrate is a subject for further study. The
orthogonality result narrows the space of possible interpretations
further: whatever the residual is, two thirds of it is the SD1/SD2
ratio and the coefficient of variation.

\section{Conclusion}

The Kuramoto residual --- the component of the phase-domain
synchronization index orthogonal to RR variability --- carries
diagnostic information beyond what conventional heart rate variability
metrics capture. At a coupling strength of $K = 2$, and across three of
five frequency distributions, the residual separates normal sinus
rhythm from atrial fibrillation with an AUC of $0.84$--$0.89$ on a
pilot corpus of eighteen paired recordings. Two of the five
distributions --- uniform random and physiological scale --- do not
exhibit the residual.

After regressing the residual on the two standard HRV features it is
most correlated with (the Poincar\'e SD1/SD2 ratio and the coefficient
of variation of the RR series), the orthogonalized residual separates
sinus from AF with AUC $0.716$, down from $0.889$. Approximately two
thirds of the residual's diagnostic value is attributable to its linear
relationship with standard HRV. One third is not.

A mechanism analysis shows why the residual is diagnostic at the
record level. Across records, $94\%$ of the between-record residual
variance in the AF cohort is explained by record-mean P-wave features
extracted from the raw ECG, features the Kuramoto engine does not have
access to. Within a single record, however, the residual's
window-to-window variation does not correlate with the same features
(mean $|r| < 0.15$ in both cohorts). The residual is therefore a
between-record discriminator, not a within-record state tracker. It
requires the full recording to stabilize (whole-record AUC 0.89;
window-level AUC 0.66--0.68 at 1--5 min; AUC 0.50 at 30 s).

The finding is preliminary. The corpus is small, the frequency
distribution dependence is not fully understood, and replication on the
full reference databases is the necessary next step. But the effect
size is large relative to the sample size, the partial-novelty result
is quantitatively characterized, and the mechanism is now understood.
The residual is a candidate feature for further investigation in
phase-domain cardiac analysis.

\begin{thebibliography}{9}

\bibitem{taskforce1996}
Task Force of the European Society of Cardiology and the North American
Society of Pacing and Electrophysiology. Heart rate variability:
standards of measurement, physiological interpretation, and clinical
use. \emph{Circulation}, 93(5):1043--1065, 1996.

\bibitem{kuramoto1975}
Y.~Kuramoto. Self-entrainment of a population of coupled non-linear
oscillators. In \emph{International Symposium on Mathematical Problems
in Theoretical Physics}, pages 420--422. Springer, 1975.

\bibitem{acebron2005}
J.~A.~Acebr\'on, L.~L.~Bonilla, C.~J.~P\'erez Vicente, F.~Ritort, and
R.~Spigler. The Kuramoto model: A simple paradigm for synchronization
phenomena. \emph{Reviews of Modern Physics}, 77(1):137--185, 2005.

\bibitem{pan1985}
J.~Pan and W.~J.~Tompkins. A real-time QRS detection algorithm.
\emph{IEEE Transactions on Biomedical Engineering},
BME-32(3):230--236, 1985.

\bibitem{brennan2001}
M.~Brennan, M.~Palaniswami, and P.~Kamen. Do existing measures of
Poincar\'e plot geometry reflect nonlinear features of heart rate
variability? \emph{IEEE Transactions on Biomedical Engineering},
48(11):1342--1347, 2001.

\bibitem{calkins2017}
H.~Calkins et al. 2017 HRS/EHRA/ECAS/APHRS/SOLAECE expert consensus
statement on catheter and surgical ablation of atrial fibrillation.
\emph{Heart Rhythm}, 14(10):e275--e444, 2017.

\bibitem{nsrdb}
R.~G.~Mark and G.~B.~Moody. MIT-BIH Normal Sinus Rhythm Database.
PhysioNet, 1990. \url{https://physionet.org/content/nsrdb/}.

\bibitem{afdb}
G.~B.~Moody and R.~G.~Mark. MIT-BIH Atrial Fibrillation Database.
PhysioNet, 1983. \url{https://physionet.org/content/afdb/}.

\bibitem{mitdb}
G.~B.~Moody and R.~G.~Mark. The impact of the MIT-BIH Arrhythmia
Database. \emph{IEEE Engineering in Medicine and Biology Magazine},
20(3):45--50, 2001.

\bibitem{richman2000}
J.~S.~Richman and J.~R.~Moorman. Physiological time-series analysis
using approximate entropy and sample entropy. \emph{American Journal of
Physiology-Heart and Circulatory Physiology}, 278(6):H2039--H2049,
2000.

\bibitem{peng1995}
C.-K.~Peng, S.~Havlin, H.~E.~Stanley, and A.~L.~Goldberger.
Quantification of scaling exponents and crossover phenomena in
nonstationary heartbeat time series. \emph{Chaos}, 5(1):82--87, 1995.

\bibitem{simoes2026engine}
A.~Simoes. The Kuramoto Engine: Physics and Mathematics of the Cardiac
Synchronization Index. Technical reference, revision 0.2, 2026.
Companion document.

\bibitem{simoes2026sensitivity}
A.~Simoes. Sensitivity Analysis Protocol: Validation Methodology for
the Kuramoto Phase Engine in NeuroCardia. Protocol, revision 0.2,
2026. Companion document.

\bibitem{simoes2026whitepaper}
A.~Simoes. NeuroCardia: An Open-Source, On-Device Phase-Domain
Diagnostic Engine for Single-Lead ECG. White paper, revision 0.2,
2026. Companion document.

\end{thebibliography}

\end{document}
